\begin{abstract}

O problema de determinar preços para um conjunto de produtos constitui um desafio central em cenários de mercados competitivos, gerando um dilema econômico: preços elevados maximizam a receita unitária, o que, consequentemente, reduz o mercado consumidor, enquanto preços baixos expandem o volume de vendas em detrimento da margem de lucro por transação. Diante disso, este trabalho estuda o \emph{Problema da Compra Mínima} (PCmin), que consiste em maximizar a receita total de um vendedor por meio da precificação de um conjunto de itens. Nesse cenário, assume-se que o valor máximo que cada comprador está disposto a pagar é previamente conhecido. Além disso, adota-se a hipótese de demanda unitária. Dessa forma, cada comprador adquire apenas o item mais barato dentre aqueles de seu interesse, ou opta por não comprar nada caso todos os preços excedam sua disposição a pagar. Embora o PCmin tenha recebido atenção crescente na literatura, sua investigação sob a perspectiva de algoritmos exatos ainda é limitada. As formulações matemáticas existentes apresentam restrições quanto à eficiência computacional, especialmente na resolução de instâncias de maior porte. Nesse contexto, identifica-se uma lacuna no desenvolvimento de métodos capazes de ampliar a escalabilidade e reduzir os tempos de solução do problema. Assim, este projeto propõe o aprimoramento de formulações de Programação Linear Inteira já existentes, bem como o desenvolvimento de abordagens heurísticas, com o objetivo de obter soluções de alta qualidade para instâncias de diferentes dimensões em tempos computacionais viáveis.


\end{abstract}

